\documentclass[uplatex,dvipdfmx]{jsarticle}
\usepackage{amsmath,amssymb}
\usepackage{bm}
\usepackage{float}
\usepackage{mathtools}
\usepackage{ascmac}
\usepackage{mathrsfs}
\usepackage{physics}
\usepackage{here}
\usepackage{url}
\usepackage{slashed}
\usepackage{tikz-cd}
\usepackage[top=25truemm,bottom=30truemm,left=20truemm,right=20truemm]{geometry}
\usepackage{xcolor}
\usepackage{titlesec}
% 文字サイズを大きくする設定
\titleformat{\section}{\huge\bfseries}{\thesection}{1em}{}
\titleformat{\subsection}{\Large\bfseries}{\thesubsection}{1em}{}
\titleformat{\subsubsection}{\Large\bfseries}{\thesubsubsection}{1em}{}
\usepackage{tcolorbox}
\tcbuselibrary{breakable, skins, theorems}
\usepackage[dvipdfmx]{hyperref}
\usepackage{pxjahyper}
\usepackage{bookmark}
\begin{document}
\title{Peskin Schroeder Chapter 12 メモ}
\author{Peskin Schroeder~~~~~~~~~~メモ作成者:~高原賢都}
\date{\today}
\maketitle
\section{Wilson's Approach to Renormalization theory}
Wilson の手法は、場の理論への汎関数積分アプローチに基づいており、そこでは量子場の自由度が積分変数となる。このアプローチでは、汎関数積分が場の短距離の自由度にどう依存するかを分離することで、紫外発散の起源を調べることができる。本節では、$\phi^4$ 理論という最も単純な例を用いてこのアイデアを説明する。\\
解析をより具体的にするため、本節では、洗練されてはいるがいくぶん謎めいた次元正則化（dimensional regularization）の手法を捨て、代わりに鋭い運動量カットオフ（sharp momentum cutoff）を用いることにする。ここでは $\phi^4$ 理論のみを扱うため、このカットオフによって Ward 恒等式を満たすことが困難になるという問題については気にしない。\\
コメント\\
($\phi^4$ 理論には、カットオフによって致命的な不整合を起こすような『ゲージ対称性』が存在しないから。)\\
Wilson の解析は、QED やこの微妙な問題が重要となる他の状況にも適用可能であるが、このアプローチの基本的な定性的結果を得るには $\phi^4$ 理論の場合で十分である。\\
9.2 節で、我々は生成汎関数 $Z[J]$ の汎関数積分表示を用いて $\phi^4$ 理論のグリーン関数を構成した。基本的な積分変数は場のフーリエ成分 $\phi(k)$ であるため、$Z[J]$ は具体的に以下の式で与えられる。
\begin{equation}
    Z[J] = \int \mathcal{D}\phi \, e^{i \int (\mathcal{L} + J\phi)} = \left( \prod_k \int d\phi(k) \right) e^{i \int (\mathcal{L} + J\phi)} \quad (12.1)
\end{equation}
鋭い紫外カットオフ $\Lambda$ を課すために、(12.1) に示された積分変数の数を制限する。すなわち、$|k| \le \Lambda$ を満たす $\phi(k)$ についてのみ積分し、$|k| > \Lambda$ では $\phi(k) = 0$ とする。汎関数積分のこのような変更は、極めて短い距離、あるいは極めて大きな運動量における量子ゆらぎの影響を評価する手法を示唆している。汎関数積分表示において、これらのゆらぎは、カットオフ付近の運動量を持つ $\phi$ のフーリエ成分についての積分によって表される。それらの変数についての積分を陽に実行してみてはどうだろうか？ そうすれば、その結果を元の汎関数積分と比較し、これらの高運動量モードが理論の物理的予言に及ぼす影響を正確に決定することができる。\\
コメント\\
(従来の繰り込み（反項を用いる方法）は、「無限大をどう取り除くか」という計算技術的な側面が強かったのに対し、このウィルソンの手法は、「エネルギー・スケールを変えたときに物理法則がどう変化するか」という物理的実体を明確に示す。この「少しずつ積分して、理論がどう変わるかを見る」というプロセスを繰り返すことで繰り込み群のアイデアを得ることができる。)\\
しかし、この解析を始める前に、一つ修正を導入しなければならない。一見すると、Minkowski 空間で紫外カットオフを定義するのが最も自然に思える。しかし、光円錐方向（lightlike directions）では、$k^2$ が小さいままでも $k$ の成分が非常に大きくなり得るため、$k^2 \le \Lambda^2$ というカットオフは大きな運動量を制御するのに完全には有効ではない。したがって、ここでは Wick 回転の後に得られるユークリッド運動量に対してカットオフが課されると考えることにする。同等のことだが、9.3 節で提示した汎関数積分のユークリッド形式を考え、その変数 $\phi(k)$（$k$ はユークリッド的）を $|k| \le \Lambda$ に制限するのである。\\
コメント\\
(例えば、$k^0 = 10^{10}$ かつ $|\vec{k}| = 10^{10}$ という非常に大きなエネルギーと運動量を持っていても、その差である $k^2$ は $0$ になる。Wick回転を行ってユークリッド空間へ移行すると、計量は $k_E^2 = (k_E^1)^2 + (k_E^2)^2 + (k_E^3)^2 + (k_E^4)^2$ となり、すべての成分が正の符号で寄与する。$k_E^2 \le \Lambda^2$ と制限すれば、すべての成分 $k_E^\mu$ は必然的に有限な範囲（$-\Lambda$ から $\Lambda$）に収まるということ。)\\
ユークリッド空間への移行はまた、第 8 章で予告した繰り込み理論と統計力学の結びつきに我々を近づける。9.3 節で見たように、$\phi^4$ 理論のユークリッド汎関数積分は、磁性体の統計力学の連続体記述と全く同じ形をしている。場 $\phi(x)$ は、ゆらぐスピン場 $s(x)$ と解釈される。実際の磁性体は原子から構成されており、原子間隔は物理的なカットオフ、すなわちゆらぎが起こり得る最短距離を提供する。カットオフされた汎関数積分は、この原子サイズの影響を粗い（crude）方法でモデル化したものになっている。\\
この類推を追求することで、場の理論における紫外カットオフの効果について、物理的な直感を得ることができる。磁性体において、原子スケールでのスピンの統計的ゆらぎを思い描くことは極めて容易である。実際、任意の臨界点から離れた温度では、統計的ゆらぎはこのスケールに制限される。数十原子間隔ほどの距離にわたって、磁性体はすでにその一様でマクロな振る舞いを示している。\\
コメント\\
(エネルギーとスケールはフーリエ変換の双対性によって結びついている。短い距離は高いエネルギーに対応し、長い距離は低いエネルギーに対応する。固体物理でのカットオフはスケールの最小である原子レベルということ。)\\
我々は第 8 章で、スピン場の相関関数をユークリッド $\phi^4$ 理論のプロパゲーターによって近似できることを見た。この近似では、
\begin{equation}
    \langle s(x)s(0) \rangle = \int \frac{d^4k}{(2\pi)^4} \frac{e^{ik \cdot x}}{k^2 + m^2} \sim_{|x| \to \infty} e^{-m|x|} \quad (12.2)
\end{equation}
温度が臨界温度から遠く離れている限り、「質量」$m$ の大きさは、この問題における唯一の自然なスケールである原子間隔によって決定される。したがって、$m \approx \Lambda$ と予想される。我々の場の理論の計算では、とりわけ $m \ll \Lambda$ となる状況に関心があり、この条件を満たすように理論のパラメータを調整した。磁性体を記述する際には、そのような調整は必要ないように思われる。\\
しかし、第 8 章で見たように、スピン場の相関が原子間隔よりもはるかに長くなり、実際に $m \ll \Lambda$ となるような状況が一つ存在する。スピン系が磁化し始める時、まさに臨界点の近傍において、ゆらぐスピンが最終的な磁化の方向を選ぼうとするにつれて、スピンは任意の長い距離にわたって相関を持つようになる。磁性体におけるこれらの長距離相関を調べるには、温度を慎重に調整し、系を相転移の近傍に持っていかなければならない。同様に、$\phi^4$ 理論のパラメータ $m$ の微調整を行い、場 $\phi(x)$ の相関が $1/\Lambda$ よりもはるかに長い距離にわたって存在するようなパラメータ領域へ量子場の理論を持っていくことを想像できる。\\
\subsection{Intergrating over a single momentum shell}
この導入を踏まえて、これから $\phi$ の高運動量自由度にわたる積分を実行する。まず、$\phi^4$ 理論の場合について、汎関数積分 (12.1) をより具体的に書き下すことから始める。\\
適切にカットオフを考えて先ほど説明した処方箋に従い、簡単のために $J=0$ と置く。すると、
\begin{equation}
    Z = \int [\mathcal{D}\phi]_\Lambda \exp\left( - \int d^d x \left[ \frac{1}{2}(\partial_\mu \phi)^2 + \frac{1}{2}m^2 \phi^2 + \frac{\lambda}{4!} \phi^4 \right] \right), \quad (12.3)
\end{equation}
ここで
\begin{equation}
    [\mathcal{D}\phi]_\Lambda = \prod_{|k| < \Lambda} d\phi(k). \quad (12.4)
\end{equation}
式(12.3)のラグランジアンにおいて、$m$ と $\lambda$ は裸のパラメータ（bare parameters）であり、したがってカウンター項（counterterms）は存在しない。表面発散度（superficial degree of divergence）を調べた時と同様に、この解析を任意の時空次元 $d$ で実行すると便利である。\\
ここで、積分変数 $\phi(k)$ を2つのグループに分割する。$b < 1$ なる割合 $b$ を選ぼう。$b\Lambda \le |k| < \Lambda$ を満たす変数 $\phi(k)$ は、これから積分を実行することになる高運動量自由度である。これらの自由度をラベル付けするために、次のように定義する。
\begin{equation}
    \hat{\phi}(k) = \begin{cases} \phi(k) & \text{for } b\Lambda \le |k| < \Lambda ; \\ 0 & \text{otherwise.} \end{cases}
\end{equation}
次に、新しい $\phi(k)$ を定義しよう。これは $|k| < b\Lambda$ に対しては古い $\phi$ と同一であり、$|k| > b\Lambda$ に対してはゼロとなるものである。すると、ラグランジアン内の古い $\phi$ を $\phi + \hat{\phi}$ に置き換えて、式(12.3)を次のように書き直すことができる。
\begin{equation}
    \begin{split}
    Z &= \int \mathcal{D}\phi \int \mathcal{D}\hat{\phi} \exp\left( - \int d^d x \left[ \frac{1}{2}(\partial_\mu \phi + \partial_\mu \hat{\phi})^2 + \frac{1}{2}m^2 (\phi + \hat{\phi})^2 + \frac{\lambda}{4!} (\phi + \hat{\phi})^4 \right] \right) \\
    &= \int \mathcal{D}\phi \, e^{-\int \mathcal{L}(\phi)} \int \mathcal{D}\hat{\phi} \exp\left( - \int d^d x \left[ \frac{1}{2}(\partial_\mu \hat{\phi})^2 + \frac{1}{2}m^2 \hat{\phi}^2 + \lambda\left( \frac{1}{6} \phi^3 \hat{\phi} + \frac{1}{4} \phi^2 \hat{\phi}^2 + \frac{1}{6} \phi \hat{\phi}^3 + \frac{1}{4!} \hat{\phi}^4 \right) \right] \right). \quad (12.5)
    \end{split}
\end{equation}
コメント\\
(Peskinでは分かりずらいが、$\phi$は新しく定義された低運動量側で値を取る場である。もともと定義されてた場を$\phi_\text{old}$とすると、$\phi_\text{old} = \phi + \hat{\phi}$ という関係が$|k| < \Lambda$ のとき成り立つ。$\hat{\phi}$は高運動量側の場である。)\\
最後の表式において、$\hat{\phi}$ に依存しないすべての項を $\mathcal{L}(\phi)$ にまとめた。異なる波長のフーリエ成分は直交するため、$\phi\hat{\phi}$ の形をした2次の項は消えることに注意されたい。\\
コメント\\
(運動量に対する制約があるので上式をフーリエ変換して積分する。4点相互作用項を運動量空間で積分する際、値を持つには4つの運動量の和がゼロになるという条件（運動量保存則）がついた積分になる。
\begin{equation}
    \int d^d x \, \phi(x)^4=\int \frac{d^d k_1 d^d k_2 d^d k_3 d^d k_4}{(2\pi)^{4d}} \delta^{(d)}(k_1+k_2+k_3+k_4) \phi(k_1) \phi(k_2) \phi(k_3) \phi(k_4)
\end{equation}
この条件では積分はキャンセルしない。しかし、2点の$\phi$ と $\hat{\phi}$ を含む項は、
\begin{equation}
    \int \frac{d^d k_1 d^d k_2}{(2\pi)^{2d}} \delta^{(d)}(k_1+k_2) \phi(k_1) \phi(k_2)
\end{equation}
となる。この積分が値を持つには、$k_1 = -k_2$ でなければならない。つまり、両方の $\phi$ が同じ運動量の大きさを持つ必要がある。しかし、$\hat{\phi}$ は $b\Lambda \le |k| < \Lambda$ の範囲で定義されているのに対し、$\phi$ は $|k| < b\Lambda$ の範囲で定義されているため、この条件を満たすことはできない。よってキャンセルする。)\\
続くいくつかの段落では、$\hat{\phi}$ についての積分をどのように実行するかを説明する。この積分により、(12.5) は次のような形の表式に変換される。
\begin{equation}
    Z = \int [\mathcal{D}\phi]_{b\Lambda} \exp\left( - \int d^d x \, \mathcal{L}_{\text{eff}} \right), \quad (12.6)
\end{equation}
ここで、$\mathcal{L}_{\text{eff}}(\phi)$ は $|k| < b\Lambda$ を満たすフーリエ成分 $\phi(k)$ のみを含む。後で見るように、$\mathcal{L}_{\text{eff}}(\phi)$ は $\mathcal{L}(\phi)$ に $\lambda$ のべき乗に比例する補正項を加えたものになる。これらの補正項は、大きな $k$ のフーリエ成分 $\hat{\phi}$ を取り除いたことを補うものであり、それまで $\hat{\phi}$ のゆらぎによって媒介されていた、残りの $\phi(k)$ 間の相互作用を供給する。\\
$\hat{\phi}(k)$ の積分を実行するために、9.2節で Feynman 則を導出する際に適用したのと同じ手法を用いる。実際、$\mathcal{L}_{\text{eff}}$ に現れる新しい項はダイアグラムの形で書けることを後で見る。この解析において、(12.5) にある $\lambda$ に比例するすべての4次項を、摂動として扱う。\\
我々は主に $m^2 \ll \Lambda^2$ の状況に関心があるため、質量項 $\frac{1}{2}m^2 \hat{\phi}^2$ も同様に摂動として扱う。すると、ラグランジアンの $\hat{\phi}$ を含む部分における主要項（leading-order term）は次のようになる.
\begin{equation}
    \int \mathcal{L}_0 = \frac{1}{2} \int_{b\Lambda \le |k| < \Lambda} \frac{d^d k}{(2\pi)^d} \hat{\phi}^*(k) k^2 \hat{\phi}(k). \quad (12.7)
\end{equation}
となる.\\
コメント\\
(運動量空間で書き換える。実空間の微分 $\partial$ は、運動量空間では $ik$ に対応するため、$(\partial \hat{\phi})^2$ は $k^2 |\hat{\phi}(k)|^2$ になる.)
\begin{equation}
    \int d^d x \left[ \frac{1}{2}(\partial \hat{\phi})^2 \right] = \frac{1}{2} \int_{b\Lambda \le |k| < \Lambda} \frac{d^d k}{(2\pi)^d} \hat{\phi}^*(k) (k^2) \hat{\phi}(k)
\end{equation}
この項は以下のプロパゲーターを導く。
\begin{equation}
    \langle \hat{\phi}(k) \hat{\phi}(p) \rangle = \frac{\int \mathcal{D}\hat{\phi} \, e^{-\int \mathcal{L}_0} \hat{\phi}(k) \hat{\phi}(p)}{\int \mathcal{D}\hat{\phi} \, e^{-\int \mathcal{L}_0}} = \frac{1}{k^2} (2\pi)^d \delta^{(d)}(k+p) \Theta(k), \quad (12.8)
\end{equation}
ここで
\begin{equation}
    \Theta(k) = \begin{cases} 1 & \text{for } b\Lambda \le |k| < \Lambda ; \\ 0 & \text{otherwise.} \end{cases}
\end{equation}
である.\\
コメント(導出)\\
分配関数は
\begin{equation}
    Z[J] = \int \mathcal{D}\hat{\phi} \, \exp\left[ -\int \mathcal{L}_0 + \int \frac{d^d q}{(2\pi)^d} J(-q)\hat{\phi}(q) \right] = \int \mathcal{D}\hat{\phi} \, \exp\left[ -S  \right]
\end{equation}
となる.~実スカラー場を仮定しているため $\hat{\phi}^*(q) = \hat{\phi}(-q)$ であることを用いる。運動量の積分範囲は $b\Lambda \le |k| < \Lambda$ であるので,~ステップ関数を用いると作用は
\begin{equation}
    S = -\frac{1}{2} \int \frac{d^d q}{(2\pi)^d} \hat{\phi}(-q) q^2 \hat{\phi}(q) \Theta(q) + \int \frac{d^d q}{(2\pi)^d} J(-q) \hat{\phi}(q) \Theta(q)
\end{equation}
となる.
\begin{equation}
    \hat{\phi}'(q) = \hat{\phi}(q) - \frac{J(q)}{q^2}
\end{equation}
という変数変換を行うと、作用は
\begin{equation}
    S = -\frac{1}{2} \int \frac{d^d q}{(2\pi)^d} \hat{\phi}'(-q) q^2 \hat{\phi}'(q) \Theta(q) + \frac{1}{2} \int \frac{d^d q}{(2\pi)^d} J(-q) \frac{1}{q^2} J(q) \Theta(q)
\end{equation}
となる.~これを分配関数に代入すると外場のない項を取り出せて
\begin{equation}
    Z[J] = Z[0] \exp\left[ \frac{1}{2} \int \frac{d^d q}{(2\pi)^d} J(-q) \frac{\Theta(q)}{q^2} J(q) \right]
\end{equation}
となる.~これを用いて、$\hat{\phi}$ の2点関数は
\begin{equation}
    \langle \hat{\phi}(k) \hat{\phi}(p) \rangle = \frac{1}{Z[0]} \frac{\delta^2 Z[J]}{\delta J(-k) \delta J(-p)}\bigg|_{J=0} 
\end{equation}
を計算する.~運動量空間の汎関数微分は
\begin{equation}
    \frac{\delta J(q)}{\delta J(p)} = (2\pi)^d \delta^{(d)}(q-p)
\end{equation}
であるので
\begin{equation}
    \frac{\delta Z[J]}{\delta J(-p)} = Z[J] \left( \frac{\Theta(p)}{p^2} J(p) \right)
\end{equation}
\begin{equation}
    \frac{\delta^2 Z[J]}{\delta J(-k) \delta J(-p)} = \left( \frac{\delta Z[J]}{\delta J(-k)} \right) \frac{\Theta(p)}{p^2} J(p) + Z[J] \frac{\Theta(p)}{p^2} \frac{\delta J(p)}{\delta J(-k)}
\end{equation}
\begin{equation}
    \left. \frac{\delta^2 Z[J]}{\delta J(-k) \delta J(-p)} \right|_{J=0} = Z[0] \frac{\Theta(p)}{p^2} (2\pi)^d \delta^{(d)}(p - (-k))
\end{equation}
よって,~
\begin{equation}
    \langle \hat{\phi}(k) \hat{\phi}(p) \rangle = \frac{1}{k^2} (2\pi)^d \delta^{(d)}(k+p) \Theta(k)
\end{equation}
が得られた.\\
式(12.5)に残っている $\hat{\phi}$ の項を摂動とみなし、指数関数を展開する。これらの摂動からの様々な寄与は、(12.8) をプロパゲーターとして用い、Wick の定理を使用することで評価できる。
\begin{equation}
    - \int d^d x \, \frac{\lambda}{4} \phi^2 \overline{\hat{\phi} \hat{\phi}} = - \frac{1}{2} \int \frac{d^d k_1}{(2\pi)^d} \mu \phi(k_1) \phi(-k_1), \quad (12.10)
\end{equation}
ここで
\begin{equation}
    \mu = \frac{\lambda}{2} \int_{b\Lambda \le |k| < \Lambda} \frac{d^d k}{(2\pi)^d} \frac{1}{k^2} = \frac{\lambda}{(4\pi)^{d/2} \Gamma(d/2)} \frac{1 - b^{d-2}}{d-2} \Lambda^{d-2}. \quad (12.11)
\end{equation}
となる.\\
コメント(行間埋め)\\
実空間の積分をフーリエ変換して運動量空間で行うと
\begin{equation}
    - \int d^d x \frac{\lambda}{4} \phi^2 \hat{\phi}^2 = -\frac{\lambda}{4} \int \frac{d^d k_1 d^d k_2 d^d p_1 d^d p_2}{(2\pi)^{4d}} (2\pi)^d \delta^{(d)}(k_1+k_2+p_1+p_2) \phi(k_1)\phi(k_2)\hat{\phi}(p_1)\hat{\phi}(p_2)
\end{equation}
となる.~Wick縮約を行うと積分は
\begin{equation}
    -\frac{\lambda}{4} \int \frac{d^d k_1 d^d k_2 d^d p_1 d^d p_2}{(2\pi)^{4d}} (2\pi)^d \delta^{(d)}(k_1+k_2+p_1+p_2) \phi(k_1)\phi(k_2) \langle \hat{\phi}(p_1) \hat{\phi}(p_2) \rangle
\end{equation}
となる.~(12.8)を代入して
\begin{equation}
    = -\frac{\lambda}{4} \int \frac{d^d k_1}{(2\pi)^d} \phi(k_1)\phi(-k_1) \int \frac{d^d p_1}{(2\pi)^d} \frac{1}{p_1^2} \Theta(p_1)
\end{equation}
となる.~ここで,~
\begin{equation}
    \mu = \frac{\lambda}{2} \int_{b\Lambda \le |k| < \Lambda} \frac{d^d k}{(2\pi)^d} \frac{1}{k^2} = \frac{\lambda}{(4\pi)^{d/2} \Gamma(d/2)} \frac{1 - b^{d-2}}{d-2} \Lambda^{d-2}
\end{equation}
と計算できるので示された.\\
(12.10) の項は、以下の指数関数の展開から生じたものとみなすこともできる。
\begin{equation}
    \exp\left( - \int d^d x \frac{1}{2}\mu \phi^2 + \dots \right). \quad (12.12)
\end{equation}
摂動級数の残りの部分も、この形に整理されることを我々はすぐに目にするだろう。したがって、係数 $\mu$ は $\mathcal{L}$ の $m^2$ 項に対する正の補正を与える。\\
コメント\\
(12.12)の1次と(12.10)右辺の項は等しいことが確認できるので,~ラグランジアンの中身を考えると
\begin{equation}
    \frac{1}{4} \lambda \phi^2 \hat{\phi}^2 \to  \frac{1}{2} \mu \phi^2
\end{equation}
となることが分かる.~よって,~$\frac{1}{4} \lambda \phi^2 \hat{\phi}^2$の1次摂動を加えることは ,~$\hat{\phi}$ の質量項$m^2$に $m^2 + \mu$ という補正(質量の繰り込み)を加えることと同等である.\\
補正項における摂動論の高次項も、同様の方法で計算できる。$\phi^4$ 理論の標準的な摂動論を導出した際と同様に、ここでもダイアグラム表記を採用すると便利である。プロパゲーター (12.8) を二重線で表そう。このプロパゲーターは、様々な4次相互作用（quartic interactions）から生じる $\hat{\phi}$ 場のペアを結びつける。これらの相互作用に現れる、積分されない場 $\phi$ を、単一の外線（external lines）として表すことにする。
$\lambda^2$ のオーダーでは、他の寄与に加えて、2つの相互作用項 $\lambda\phi^2\hat{\phi}^2$ の縮約を含む項が現れる。各項は、2本の実線と2本の二重線を結ぶ頂点に対応する。可能な縮約は2通りある。\\
これらのうち、最初の非連結図形（disconnected diagram）は、指数関数 (12.12) における $\lambda^2$ オーダーの項を与える。2つ目は新しい寄与であり、$\mathcal{L}(\phi)$ における $\phi^4$ 相互作用への補正となる。\\
では、この2つ目の寄与を評価してみよう。簡単のため、$\phi$ の因子が運ぶ外線運動量が $b\Lambda$ に比べて非常に小さく、無視できる極限を考える。すると、この図形は次の値を持つ.
\begin{equation}
    -\frac{1}{4!} \int d^dx \zeta \phi^4 \qquad (12.14)
\end{equation}
コメント\\
(これは,~12.13の右のダイアグラムを計算した結果になる。12.10では縮約1つ分の寄与を考えたが、この積分では2つの縮約と$\phi^4$の積分なので、対応する項は$\phi^4$の項であり相互作用$\lambda$は縮約2つ分の計算をすることになる。この計算を12.10と同様に行うと)
\begin{equation}
    \zeta = -i \frac{\lambda^2}{2} \left(\frac{4!}{2!2!} \right)^2 \int_{b\Lambda \le |k| < \Lambda} \frac{d^dk}{(2\pi)^d} \left( \frac{i}{k^2} \right)^2 \approx -\frac{3\lambda^2}{(4\pi)^{d/2}\Gamma(d/2)} \frac{(1-b^{d-4})}{d-4} \Lambda^{d-4} \xrightarrow{d\to4} -\frac{3\lambda^2}{16\pi^2} \log \frac{1}{b} \qquad (12.15)
\end{equation}
となる。分子の2は可能な2つの縮約を数えたものである,外線や頂点を数えることによる追加の組み合わせ因子はない。10.2節の $\phi^4$ 理論の解析において、運動量 0 から $\Lambda$ の範囲で積分された似たような図形に遭遇し、それは対数的な紫外発散を生じた。Wilson の扱いにおいて、この発散は病的なものではなく、単にその図形がすべての運動量スケールから寄与を受けていることの現れである。実際、それは運動量スケール $m$ と $\Lambda$ の間の各対数区間から等しい寄与を受ける。後で見るように、各運動量区間からこの図形への（有限の）寄与は自然な物理的重要性を有している。\\
ここまで説明した図形を用いた摂動論は、$\phi^2$ や $\phi^4$ に比例する寄与だけでなく、$\phi$ のより高次のべき乗への寄与も生成する。例えば、次の図形は $\phi^6$ 相互作用への寄与を生成する.
また、図形の外線運動量をもはや無視しなくなったときに生じる微分子相互作用もある。より厳密な扱いでは、これらの運動量でテイラー展開することになる；例えば、式(12.14)に加えて、外線運動量の2乗に比例する項が得られ、それは次のように書き直すことができる.
\begin{equation}
    -\frac{1}{4} \int d^dx \, \eta \phi^2 (\partial_\mu \phi)^2 \qquad (12.17)
\end{equation}
また、$\phi$ が運ぶ運動量の4乗、6乗、およびそれ以上の項も見つかるだろう。一般に、$\hat{\phi}$ を積分消去する手続きは、場 $\phi$ とその微分のあらゆる可能な相互作用を生成する。\\
図形の寄与は、それらを指数関数として再総和（resumming）することでわずかに単純化できる。(12.13)ですでに見られたように、我々の図形展開は非連結図形を生成する。式(4.52)で用いたのと同じ組み合わせの議論により、級数の和を連結図形（connected diagrams）の和の指数関数として書き直すことができる。これはまさに式(12.6)を導き、以下のようになる。
\begin{equation}
    \mathcal{L}_{\text{eff}} = \frac{1}{2}(\partial_\mu \phi)^2 + \frac{1}{2}m^2 \phi^2 + \frac{1}{4!}\lambda \phi^4 + (\text{連結図形の和}) \qquad (12.18)
\end{equation}
図形の寄与には、$m^2$ と $\lambda$ への補正に加えて、あらゆる可能な高次元演算子（higher-dimension operators）が含まれる。これで、新しいラグランジアン $\mathcal{L}_{\text{eff}}(\phi)$ を使って $\phi(k)$ の相関関数を計算したり、S行列要素を計算したりすることができる。$\phi(k)$ は $b\Lambda$ までの運動量しか含んでいないため、そのような計算におけるループ図形は、その引き下げられたカットオフまでしか積分されない。(12.18)の補正項は、まさにこの変化を補償するものである。\\
式(12.18)における高次元演算子の出現に戸惑うかもしれない。我々は、$\phi^4$ 理論の元のラグランジアンが、繰り込み可能な相互作用（renormalizable interactions）のみを含むように選んだ。一見したところ、変数 $\hat{\phi}$ を積分消去したときに、あらゆる可能な繰り込み不可能な相互作用が現れるのは気がかりである。しかし、後で見るように、我々の手続きは実際にはこれらの繰り込み不可能な相互作用の寄与をコントロール下に保つ。事実、我々の解析は、非常に大きなカットオフ $\Lambda$ と共に用いると定義された元のラグランジアンに繰り込み不可能な相互作用が存在していても、$\Lambda$ よりもはるかに小さなスケールの物理には無視できるほどの影響しか与えないことを意味する。
\subsection{Renormalization Group Flows}
ここで、新しい汎関数積分(12.6)と、我々が出発点とした汎関数積分(12.3)をより慎重に比較してみよう。これを行う最も便利な方法は、(12.6)における距離と運動量を次のようにリスケールすることである。
\begin{equation}
    k' = k/b, \qquad x' = xb, \qquad (12.19)
\end{equation}
そうすることで、変数 $k'$ は $|k'| < \Lambda$ にわたって積分されることになる。\\
式(12.18)の形は、図式的に次のように書ける。
\begin{equation}
    \int d^d x \mathcal{L}_{\text{eff}} = \int d^d x \left[ \frac{1}{2}(1 + \Delta Z)(\partial_\mu \phi)^2 + \frac{1}{2}(m^2 + \Delta m^2)\phi^2 + \frac{1}{4!}(\lambda + \Delta \lambda)\phi^4 + \Delta C (\partial_\mu \phi)^4 + \Delta D \phi^6 + \dots \right] 
\end{equation}
再スケーリングされた変数 $x'$ を用いると、これは次のようになる.
\begin{equation*}
    \int d^d x' b^{-d} \left[ \frac{1}{2}(1 + \Delta Z)b^2 (\partial'_\mu \phi)^2 + \frac{1}{2}(m^2 + \Delta m^2)\phi^2 + \frac{1}{4!}(\lambda + \Delta \lambda)\phi^4 + \Delta C b^4 (\partial'_\mu \phi)^4 + \Delta D \phi^6 + \dots \right]
\end{equation*}
この解析を通じて、我々は最初の項（運動項）以外のすべての項を小さな摂動として扱ってきた。元の結合定数が小さい限り、これは(12.21)を扱う上でも妥当な近似である。\\
式(12.21)の元の汎関数積分は、プロパゲータ(12.8)を導いた。もし、スカラー場 $\phi$ を次のように再スケーリングすれば、新しい作用(12.21)は全く同じプロパゲータを与えることになる.
\begin{equation}
    \phi' = [b^{2-d}(1 + \Delta Z)]^{1/2} \phi \quad (12.22)
\end{equation}
この再スケーリングの後、非摂動（自由場）の作用は元の形式に戻り、一方で様々な摂動項は変換を受ける。
\begin{equation}
    \int d^d x \mathcal{L}_{\text{eff}} = \int d^d x' \left[ \frac{1}{2}(\partial'_\mu \phi')^2 + \frac{1}{2} m'^2 \phi'^2 + \frac{1}{4!} \lambda' \phi'^4 + C' (\partial'_\mu \phi')^4 + D' \phi'^6 + \dots \right] \quad (12.23)
\end{equation}
ラグランジアンの新しいパラメータは以下の通りである.
\begin{equation}
    m'^2  = (m^2 +\Delta m^2)(1+\Delta Z)^{-1} b^{-2}
\end{equation}
\begin{equation}
    \lambda' = (\lambda + \Delta \lambda)(1+\Delta Z)^{-2} b^{d-4}
\end{equation}
\begin{equation}
    C' =( C+ \Delta C )(1+\Delta Z)^{-2} b^{d}
\end{equation}
\begin{equation}
    D' = (D + \Delta D)(1+\Delta Z)^{-3} b^{2d-6}
\end{equation}
などなど。（元のラグランジアンは $C=D=0$ であったが、初期値の $C$ や $D$ がゼロでなくても、同じ方程式が適用される。）\\
コメント\\
($\Delta$が付いてる量は$\hat{\phi}$の摂動、つまり積分によってパラメータに繰り込まれた量であり、$\Delta$が付いてない量は元のラグランジアンのパラメータである。)\\
すべての補正項 $\Delta m^2, \Delta \lambda$ 等は、ファインマン図形から生じるため、摂動論が正当化される範囲では小さな量である。\\
運動量空間の「殻（シェル）」を積分して取り除く操作と、(12.19)の再スケーリングを組み合わせることで、我々はこの操作をラグランジアンの変換として書き換えた。この手順を繰り返すことで、さらに別の運動量空間の殻を積分し、ラグランジアンをさらに変換させることができる。\\
連続的な積分を繰り返すことで、(12.24)の変換のさらなる反復が生じる。パラメータ $b$ を 1 に非常に近く設定すれば、積分される殻は無限に薄くなり、変換は連続的なものとなる。すると、高運動量の自由度を積分した結果を、すべての可能なラグランジアンの空間における「軌跡」または「流れ（フロー）」として記述できるようになる。\\
コメント\\
($b$を1に近づけることでシェルは薄くなり、変換式を漸化式のように再帰的に適用すると運動量を段階的に下げることができる。$b$は1よりも小さいため、各パラメータの挙動はオーダーで評価できる。)\\
歴史的な理由から、これらの連続的に生成されるラグランジアンの変換は、繰り込み群（Renormalization Group）と呼ばれている。自由度を積分する操作は可逆ではないため、厳密な意味での「群」ではない。一方で、第13章で見るように、それらは確かに半群（semigroup）を形成している。\\
外線運動量 $p_i$ がすべて $\Lambda$ よりはるかに小さい場の相関関数を計算したいとしよう。この相関関数は、元のラグランジアン $\mathcal{L}$ を用いて全運動量について積分して計算することもできるし、有効ラグランジアン $\mathcal{L}_{\text{eff}}$ を用いて $p < \Lambda/b$ の運動量について積分して計算することもできる。どちらの手順も最終的には同じ結果を与えるはずである。\\
しかし、前者の手順では、高運動量のゆらぎの効果はループ図形を計算するまで現れない。後者の手順では、これらの効果はすでに新しい結合定数（$m', \lambda', C', \dots$）に吸収されている。したがって、それらの影響はラグランジアンから直接読み取ることができる。\\
最も単純なケースは、点 $m^2 = \lambda = C = D = \dots = 0$ の近傍にあるラグランジアンを考えることである。ここではすべての摂動が消失する。我々は、この点が不変に保たれるように変換を定義した。この点を自由場ラグランジアンの固定点と呼ぶ。\\
\begin{equation}
    \mathcal{L}_0 = \frac{1}{2}(\partial_\mu \phi)^2 \quad (12.25)
\end{equation}
原点でのラグランジアンは、すべての高次元演算子が消えるため、単純な自由場のラグランジアンになる。$\mathcal{L}_0$ の近傍では、(12.24)の反復方程式において $\Delta m^2, \Delta \lambda$ などの項を無視し、摂動に対して線形な項のみを残すことができる。これにより、特に単純な変換則が得られる。
\begin{equation}
    m'^2 = m^2 b^{-2}, \qquad \lambda' = \lambda b^{d-4}, \qquad C' = C b^d, \qquad D' = D b^{2d-6}, \dots \quad (12.26)
\end{equation}
$b < 1$ であるため、$b$ の負のべき乗が掛けられるパラメータは増大し、一方で $b$ の正のべき乗が掛けられるパラメータは減少する。もしラグランジアンに増大する係数が含まれていれば、それらは最終的に系を $\mathcal{L}_0$ から遠ざけることになる。\\
有効ラグランジアンにおける様々な項を、 $\mathcal{L}_0$ への摂動として加えられる局所演算子の集合として語るのが慣習である。反復手順の間に係数が増大する演算子を重要演算子（relevant operators）、係数が減少する（消えていく）演算子を不要演算子（irrelevant operators）と呼ぶ。\\
例えば、スカラー場の質量演算子 $\phi^2$ は常に「重要」であり、一方で $\phi^4$ 演算子は $d < 4$ の場合に「重要」となる。ある演算子の係数が $b^0$ 倍される場合、その演算子は周辺的（marginal）であると呼ばれる。例えば、$\phi^4$ 演算子は $d=4$ の場合に周辺的である。\\
その係数が増大するか減少するかを知るには、高次の補正の影響を含める必要がある。一般に、 $N$ 個の $\phi$ と $M$ 個の微分を持つ演算子の係数は、次のように変換される。
\begin{equation}
    C'_{N,M} = b^{N(d/2 - 1) + M - d} C_{N,M} \quad (12.27)
\end{equation}
導出\\
\begin{equation}
    \int d^d x \, C_{N,M} \, (\partial^M \phi^N)
\end{equation}
を再スケーリングされた変数 $x'$ を用いて書き換えると、ラグランジアンの運動項を元の形式に保つための $\phi$ の再スケーリングを考慮して
\begin{equation}
    \phi = b^{(d-2)/2} \phi'
\end{equation}
となるので、上の式は
\begin{equation}
    \int d^d x' b^{-d} C_{N,M} b^{N(d-2)/2} b^M (\partial'^M \phi'^N) = \int d^d x' b^{N(d/2 - 1) + M - d} C_{N,M} (\partial'^M \phi'^N)
\end{equation}
となる.\\
この係数はちょうど $(d - d_{N,M})$ であることに注意せよ。ここで $d_{N,M}$ は第10.1節の最後で計算した演算子の質量次元である。言い換えれば、「重要」「不要」「周辺的」な演算子は、自由理論 $\mathcal{L}_0$ の近傍では、第10.1節のべき乗計数による解析における「繰り込み可能（renormalizable）」「非繰り込み可能（non-renormalizable）」「超繰り込み可能（super-renormalizable）」な相互作用項に正確に対応している。\\
係数の変化は、自由場固定点の近傍での単純な次元解析を用いても理解できる。質量次元 $d_i$ を持つ係数を持つ演算子は、(質量)$^{d-d_i}$ という次元の係数を持つ。この質量の自然な大きさの尺度はカットオフ $\Lambda$ である。したがって、$d_i < d$ であれば、その摂動の相対的な大きさは低運動量においてますます重要になる。一方、$d_i > d$ であれば、運動量 $p \to 0$ においてこの項の相対的なサイズは $(p/\Lambda)^{d_i - d}$ として減少する。この意味で、そのような項は真に「不要」なのである。\\
我々は今、少なくともゼロ結合固定点の近傍では、カットオフのスケールにおいて任意に複雑なラグランジアンであっても、わずかな数の「重要」な相互作用のみを含むラグランジアンへと退化することを示した。この結果を第10章の結論と比較するのは興味深い。そこでは、$\Lambda \to \infty$ という極限をできるだけ早く取るべきだという哲学をとった。その極限が、負の質量次元を持つパラメータを含まないラグランジアンに対してのみ、うまく定義された（well-defined）予測を与えることを見出した。この観点からすると、例えばQEDにそのようなパラメータが含まれていないのは、そうでなければ理論がうまく定義された予測を与えないため、非常に幸運なことであるように思われた。ウィルソンの解析は、まさに逆の視点をとる。すなわち、いかなる量子場理論も、物理的な意味を持つカットオフ $\Lambda$ を用いて根本的に定義されるという視点である。統計力学の応用では、この運動量スケールは逆格子間隔である。QEDや他の素粒子物理学の量子場理論では、カットオフは時空の何らかの根本的な粒状性（恐らく量子重力のゆらぎの結果）に関連しているはずである。\\
しかし、このスケールが何であれ、それは現在の実験の到達範囲をはるかに超えている。我々が今示した議論は、QEDや他の量子場理論の「繰り込み可能性」が、それらの物理的相互作用の本質的なスケールを説明していることを示している。QEDの基礎となるラグランジアンが、その結合定数が十分に弱ければ何であれ、我々の実験のエネルギーにおいては、繰り込み可能な有効ラグランジアンとして記述されなければならないのである。\\
一方、これらの単純な結論は、十分に強い場理論の相互作用によって変化し得ることも強調しておくべきである。自由場固定点から離れると、単純な変換則(12.26)は、結合定数の高次べきに比例する補正を受ける。これらの補正が十分に大きければ、繰り込み群の流れを止めたり逆転させたりすることさえあり得る。それらは新しいタイプの $\Lambda \to \infty$ 極限を与えるような、新しい種類の固定点を作り出すことさえあるだろう。\\
比較的単純な文脈で相互作用の及ぼし得る影響を例示するために、特定の $\phi^4$ 理論について、$\mathcal{L}_0$ 近傍の繰り込み群の流れを議論しよう。$d > 4$、$d = 4$、$d < 4$ の3つのケースを順に考えるのが有益である。\\
$d > 4$ のとき、唯一の重要な演算子はスカラー場質量項である。そのとき、繰り込み群の流れは図12.1(a)のように $m^2$ 軸の近くを通る。$\phi^4$ 相互作用（および起こり得る高次の相互作用）は消え去り、一方で質量項は重要性を増していく。\\
以前の章では、常に質量がカットオフに比べて小さい極限での $\phi^4$ 理論を議論してきた。この条件を繰り込み群の流れの言葉で書き換えてみよう。流れの過程で、有効質量 $m'^2$ は増大し、最終的には現在のカットオフに等しくなる。例えば、自由場固定点の近傍では、 $n$ 回の反復の後、 $m'^2 = m^2 b^{-2n}$ となり、最終的には $m'^2 \sim \Lambda^2$ となる $n$ が存在する。この時点で、元の $\Lambda$ から有効質量のスケールの間の全運動量領域を積分し終えたことになる。質量項は、それ以降の残りの量子ゆらぎを抑制する。一般に、スカラー場の質量がカットオフに比べて小さいという基準は、非常に多数の反復の後でのみ $m'^2 \sim \Lambda^2$ となるという言明と同等である。\\
この基準は、繰り込み群のフローの初期条件が調整され、軌跡が固定点のすぐ近くを通る場合には常に満たされる。原理的には、フローはレリバントでない（irrelevantな）演算子の方向に沿って、遠く離れたところから始まる可能性がある。元の $m^2$ の値は、$\mathcal{L}_{\text{eff}}$ へのファインマン・ダイアグラムからの寄与による補正によってキャンセルされる限り、特に小さくなくてもよい。したがって、$d > 4$ において複雑な非線形ラグランジアンを書き、ただし元の $m^2$ を調整して、このラグランジアンから始まる軌跡が最終的に自由場の固定点 $\mathcal{L}_0$ の近くを通るようにすることで、スカラー場の理論を構築することが想像できる。この場合、カットオフに比べて運動量が小さいところでの有効理論は極めて単純なものになるはずである。つまり、無視できるほどの非線形相互作用しか持たない自由場の理論になるだろう。次章で議論するように、この驚くべき予測は、4次元より大きい次元の磁性体の数理モデルにおいて検証されている。元のモデルが極めて非線形であるにもかかわらず、相転移付近のスピンの相関関数は、式 (12.2) の高次元における類似物によって与えられる自由場の形を持つのである。\\
次に $d = 4$ の場合を考えてみよう。この場合、式 (12.26) は、長距離において $\phi^4$ 相互作用が重要かそうでないかを判断するのに十分な情報を与えてくれない。したがって、完全な変換則 (12.24) に戻らなければならない。$\Delta\lambda$ への主要な寄与は式 (12.15) で与えられる。$\Delta Z$ への主要な寄与は $\lambda^2$ のオーダーであり、無視することができる。（これはまさにセクション 10.2 の $\delta z$ に対する最初の補正で起こったことである。）したがって、以下の変換が得られる。
\begin{equation}
    \lambda' = \lambda - \frac{3\lambda^2}{16\pi^2} \log(1/b). \quad (12.28)
\end{equation}
これは、高運動量の自由度を積分消去していくと、$\lambda$ がゆっくりと減少することを意味している。\\
補正 $\Delta\lambda$ に寄与するダイアグラムは、セクション 10.2 で計算した1ループダイアグラムと同じ構造を持っている。実際、これらは本質的に同じダイアグラムであり、積分を反復的に実行するか一度に実行するかという点のみが異なる。しかし、セクション 10.2 のダイアグラムが紫外発散を持っていたのに対し、ウィルソンのアプローチにおける対応するダイアグラムは well-defined（きちんと定義されている）であり、結合定数の単純な発展方程式の係数を与える。この変換は、本章で行う紫外発散の再解釈の最初の例となる。\\
変換則 (12.28) は、$\mathcal{L}_0$ 付近の繰り込み群フローが図 12.1(b) に示すような、ゆっくりと減衰する1つの方向を持つ形になることを意味している。フローを十分に遠くまでたどると、その振る舞いは再び自由場のようになるはずである。この描像は、カットオフを無限大にする極限において、4次元の相互作用する $\phi^4$ 理論が存在しないという不可解な意味合いを持っている。この結果についてさらに議論し、それにもかかわらず $\phi^4$ 理論をモデル場の理論として使用することが妥当である理由を、セクション 12.3 で説明する。\\
最後に$d < 4$ の場合を考えてみよう。この場合、$\lambda$はrelevantである。\\
したがって、自由度を積分消去していくと、理論は自由理論 $\mathcal{L}_0$ から離れるように流れていく。つまり長距離においては、$\phi^4$ 相互作用がますます重要になる。しかし、$\lambda$ が大きくなると、式 (12.28) に示されるような非線形な補正も考慮しなければならない。$d < 4$ においてこの特定の効果を含めると、以下の漸化式が得られる.
\begin{equation}
    \lambda' = \left[ \lambda - \frac{3\lambda^2}{(4\pi)^d \Gamma(d/2)} \frac{b^{d-4}-1}{4-d} \right] b^{4-d}. \quad (12.29)
\end{equation}
この方程式は、リスケーリングによる増加が非線形効果による減少によって相殺されるような $\lambda$ の値が存在することを意味している。この値において、自由度を積分消去しても $\lambda$ は変化しない。対応するラグランジアンは、繰り込み群フローの第2の固定点である。極限 $d \rightarrow 4$ において、フロー (12.29) は (12.28) に近づき、そのため新しい固定点は自由場の固定点と融合する。$d$ が十分に 4 に近い場合、新しい固定点は、反復によって質量パラメータ $m^2$ が増加するという性質を $\mathcal{L}_0$ と共有するだろう。その場合、質量演算子は新しい固定点付近でレリバントな演算子となり、その結果、繰り込み群フローは図 12.2 に示すような形になる。
この例では、繰り込み群の新しい固定点は、変換方程式が摂動論で計算できるほど十分に弱い結合を持つラグランジアンを有していた。原理的には、ラグランジアンが強結合であり、そのため繰り込み群変換がファインマン・ダイアグラムによる解析では理解できないような固定点を見つけることもできる。そのような固定点の多くの例は、2次元の厳密に解けるモデル場の理論において知られている。$^\dagger$ しかし、現在に至るまで、物理的応用において重要な量子場の理論の例はすべて、自由場の固定点によって支配されているか、あるいは前の段落で記述したような、特定の極限において自由場の固定点に近づく固定点によって支配されていることが判明している。なぜこうなるのかは誰にも分かっていない。
\section{The Callan-Symanzik Equation}
\subsection{The Callan-Symanzik Equation}



\end{document}